\subsection{Scales and independent angles}\label{p:sec:scales}

We choose the finite polynomial accuracy before choosing each new block.
This leaves enough freedom to obtain both the sampling estimates and the
threshold crossings needed by Lemma~\ref{lem:conditional-convolution}.
All choices in this subsection are independent of any invariant projection.

Partition $\N=\{1,2,\ldots\}$ into consecutive finite blocks $J_j$ of
length $m_j$, writing
\[
 d_0=0,\qquad d_j=\sum_{s\le j}m_s,\qquad
 J_j=\{d_{j-1}+1,\ldots,d_j\}.
\]
The cutoff $K_j$ will specify the finite coefficient window in the
contradiction, while $D_j$ and $\varepsilon_j$ will give the coefficient
bound and the product lower bound from Lemma~\ref{lem:polynomial-coefficient}.  After these
are fixed, $L_j$ specifies the required decay of the omitted terms;
the logarithmic drop $C_j$ and the block length $m_j$ are chosen to
achieve that decay while retaining uniform bounds on the coupling
marginals.

We now choose these parameters.  Put $\kappa_j=1/(100j)$ and
start with $K_1=100$ and $G_1=1$.  At stage $j$, the numbers $K_j,G_j$
and all parameters of earlier blocks have already been fixed.  Make
the following choices in the stated order.

\begin{enumerate}
\item Define
\[
 D_j=\exp\left((K_j+1)\sum_{s<j}C_sm_s\right).
\]
Apply Lemma~\ref{lem:polynomial-coefficient} with $K=K_j$, $D=D_j$, and
$\delta=1/j$ to choose $\varepsilon_j>0$.  This is possible because
$D_j\ge1\ge1/j$.  Consequently,
\begin{equation}\label{p:eq:poly-lower}
 \max_{0\le n\le2K_j}|(AB)_{-n}|\ge\varepsilon_j
\end{equation}
whenever $A,B$ are supported in $[-K_j,K_j]$, their coefficients
have modulus at most $D_j$, and $|A_0|,|B_0|\ge1/j$.

\item Choose $L_j\ge1$, increasing with $j$, so large that
\begin{equation}\label{p:eq:epsilon-choice}
 (2K_j+1)3\exp(-L_j/10000)\le\frac{\varepsilon_j}{j}.
\end{equation}
Set $C_j=\max\{L_j,K_j^2\}$, and choose a finite $G_{j+1}\ge1$
such that
\begin{equation}\label{p:eq:G-choice}
 \kappa_{j+1}\sqrt{G_{j+1}}\ge L_j.
\end{equation}

\item Choose an integer $m=m_j$ sufficiently large, and set
\[
 K_{j+1}=\left\lceil m^{(3/4)m}\right\rceil,
\]
subject to
\begin{equation}\label{p:eq:scale-conditions}
 \begin{gathered}
 m\ge\max\{K_j,2^j,4G_j,100\},\qquad d_{j-1}\le m/100,\\
 \kappa_jm^{1/40}\ge L_j,\qquad
 G_{j+1}K_{j+1}<m^{(4/5)m}.
 \end{gathered}
\end{equation}
\end{enumerate}

Every choice is finite and feasible.  Indeed, the first two steps do
not depend on the new $m_j$.  All but the final requirement in
\eqref{p:eq:scale-conditions} are lower bounds on $m$.  For the final
one, the already fixed number $G_{j+1}$ satisfies
\[
 \frac{G_{j+1}\lceil m^{(3/4)m}\rceil}{m^{(4/5)m}}
 \le 2G_{j+1}m^{-m/20}\longrightarrow0.
\]
There is therefore no upper restriction conflicting with the lower
bounds.

Write
\begin{equation}\label{p:eq:N-and-widths}
 N_j=m_jK_j,\qquad a_j=G_j/m_j,\qquad b_j=m_j^{-19/20}.
\end{equation}
The construction gives the useful consequences
\begin{equation}\label{p:eq:scale-consequences}
 \begin{gathered}
 m_j\ge2^j,\qquad d_j\le1.01m_j,\qquad N_j\le m_j^2,\\
 K_{j+1}>2K_j,\qquad
 \sum_{s>j}\frac1{K_s}\le\frac2{K_{j+1}},\qquad
 \frac1{m_j}\le a_j,b_j\le\frac14.
 \end{gathered}
\end{equation}
For example, $m_j\ge100$ and $m_j\ge K_j$ imply
$K_{j+1}\ge m_j^{(3/4)m_j}>2m_j\ge2K_j$, and the reciprocal
estimate follows by summing a geometric majorant.  Also
$m_{j+1}\ge K_{j+1}>m_j$, so each of $m_j,K_j,N_j$ is increasing.
The upper bound for $a_j$ follows from $m_j\ge4G_j$; the other width
bounds follow from $G_j\ge1$ and $m_j\ge100$.

Use the coordinate order on $X$ defined above.

On a separate product probability space choose independent random
angles $\theta_i$, with $\theta_i$ uniform on $[1/N_j,2/N_j]$ when
$i\in J_j$, and let
\[
 \sigma_\theta x=x+\theta\quad\text{in }X.
\]
For every realization,
\begin{equation}\label{p:eq:angle-summability}
 \sum_{i=1}^\infty\theta_i
 \le2\sum_{j=1}^\infty\frac{m_j}{N_j}
 =2\sum_{j=1}^\infty\frac1{K_j}<\infty.
\end{equation}
Almost surely $1,\theta_1,\theta_2,\ldots$ are rationally independent.
To see this, multiply any proposed finite rational relation by a
common denominator.  If some angle has a nonzero integer coefficient,
conditioning on all other angles leaves one prescribed value of an
atomless random variable, an event of probability zero.  A relation
with no angle term cannot hold unless its constant term is zero.
There are only countably many nontrivial finite integer relations.

For a rationally independent realization the rotation is ergodic
and free.  Indeed, the characters
\[
 x\longmapsto\exp\left(2\pi i\sum_i v_ix_i\right),
 \qquad v_i\in\Z,\quad v_i=0\text{ for all but finitely many }i,
\]
form an orthonormal basis of $L^2(\mu)$: finite-coordinate functions
are dense, and the usual Fourier bases are complete on each finite
torus.  Translation multiplies this character by
$\exp(2\pi i\sum_i v_i\theta_i)$, which equals $1$ only for the
constant character.  Hence the invariant $L^2$ functions are constant,
which is ergodicity.  Also $n\theta_1\notin\Z$ for every
$n\in\Z\setminus\{0\}$, so $\sigma_\theta^n x\ne x$ for every
$x\in X$.

\begin{proposition}[Phase ranges]\label{p:prop:phases}
Almost surely, for every sufficiently large $j$ the following hold
simultaneously for all integers $k$ in the indicated ranges, with
$m=m_j$:
\begin{enumerate}
\item If $G_jK_j\le |k|\le N_j/4$, then every $i\in J_j$ satisfies
\[
 1-\{k\theta_i\}\in[a_j,1-a_j].
\]
\item If $N_j/4\le |k|\le m^{(4/5)m}$, then at least $m/20$
coordinates $i\in J_j$ satisfy
\[
 1-\{k\theta_i\}\in[b_j,1-b_j].
\]
\end{enumerate}
The first assertion in fact holds at every stage and for every
realization of the angles in their prescribed intervals.
\end{proposition}

\begin{proof}
In the first range, put $t=|k|\theta_i$.  The angle bounds give
\[
 a_j=G_j/m_j\le t\le1/2.
\]
For $k>0$ the threshold is $1-t$, and for $k<0$ it is $t$.
Both belong to $[a_j,1-a_j]$ because $a_j\le1/4$.

For the second range, the uniform law of $k\theta_i$ on an interval
of length $|k|/N_j$ has, after reduction modulo one, a density bounded
by
\[
 1+N_j/|k|\le5.
\]
This follows by counting the integer translates covering a point:
the covering multiplicity of an interval of length $r$ is at most
$r+1$.  The statement also applies to negative $k$ by reflection.
Thus a given coordinate misses $[b_j,1-b_j]$ with probability at
most $10b_j$.  If fewer than $m/20$ coordinates satisfy the assertion,
more than $19m/20$ coordinates miss the interval.  Independence over
$i\in J_j$ and a union bound over the possible sets of misses show
that, for a fixed $k$, this failure probability is at most
$2^m(10b_j)^{(19/20)m}$ for large $m$.  A further union bound over
the at most $2m^{(4/5)m}$ nonzero integers in the range gives
\[
 \mathbb P(\text{failure at stage }j)
 \le2m^{(4/5)m}2^m(10b_j)^{(19/20)m}
 =\exp\left(-\frac{41}{400}m\log m+O(m)\right).
\]
These probabilities are summable because $m_j\ge2^j$.
The first Borel--Cantelli Lemma now gives the second assertion for
all sufficiently large $j$, simultaneously for every integer in
each stated range.
\end{proof}

We retain the random angles until the sampling construction below
has supplied one further probability-one property.
